% =====================================================================
%  Monster Tamers — Architettura del backend e decisioni di progetto
%  Advanced Software Engineering (Lab) — Università di Pisa, A.A. 2026/27
%  Compilazione: tectonic -X compile architettura_backend.tex
% =====================================================================
\documentclass[11pt,a4paper]{article}

% ---------------------------------------------------------------------
%  Pacchetti
% ---------------------------------------------------------------------
\usepackage{fontspec}
\usepackage[italian]{babel}
\usepackage[a4paper,top=2.6cm,bottom=2.6cm,left=2.5cm,right=2.5cm,headheight=14pt]{geometry}
\usepackage{microtype}
\usepackage{amsmath}
\usepackage{float}
\usepackage{xcolor}
\usepackage{array,booktabs,tabularx,xltabular}
\usepackage{enumitem}
\usepackage{titlesec}
\usepackage{fancyhdr}
\usepackage[font=small,labelfont=bf,labelsep=period]{caption}
\usepackage{adjustbox}
\usepackage{needspace}
\usepackage{tikz}
\usetikzlibrary{arrows.meta,positioning,shapes.geometric,fit,calc}
\usepackage{url}
\usepackage[hidelinks]{hyperref}

% ---------------------------------------------------------------------
%  Font
% ---------------------------------------------------------------------
\setmainfont{texgyrepagella}[
  Extension      = .otf,
  UprightFont    = *-regular,
  BoldFont       = *-bold,
  ItalicFont     = *-italic,
  BoldItalicFont = *-bolditalic]
\setsansfont{texgyreheros}[
  Extension      = .otf,
  UprightFont    = *-regular,
  BoldFont       = *-bold,
  ItalicFont     = *-italic,
  BoldItalicFont = *-bolditalic]
\setmonofont{lmmono10-regular}[Extension=.otf, Scale=MatchLowercase]
\usepackage{unicode-math}
\setmathfont{texgyrepagella-math.otf}

% ---------------------------------------------------------------------
%  Impaginazione
% ---------------------------------------------------------------------
\definecolor{rulegray}{gray}{0.45}

\titleformat{\section}{\sffamily\Large\bfseries}{\thesection}{0.8em}{}
\titleformat{\subsection}{\sffamily\large\bfseries}{\thesubsection}{0.8em}{}
\titleformat{\subsubsection}{\sffamily\normalsize\bfseries}{\thesubsubsection}{0.8em}{}
\titlespacing*{\section}{0pt}{2.2ex plus 1ex minus .2ex}{1.2ex plus .2ex}
\titlespacing*{\subsection}{0pt}{1.8ex plus 1ex minus .2ex}{0.8ex plus .2ex}

\pagestyle{fancy}
\fancyhf{}
\fancyhead[L]{\small\sffamily Monster Tamers}
\fancyhead[R]{\small\sffamily Architettura del backend}
\fancyfoot[C]{\small\sffamily\thepage}
\renewcommand{\headrulewidth}{0.4pt}

\setlength{\parindent}{0pt}
\setlength{\emergencystretch}{3em}
\setlength{\parskip}{0.55em}
\setlist{itemsep=0.15em, topsep=0.3em}
\renewcommand{\arraystretch}{1.25}

\newcolumntype{L}[1]{>{\raggedright\arraybackslash}p{#1}}
\newcolumntype{Y}{>{\raggedright\arraybackslash}X}

% Identificatori, campi, valori
% (spezzano la riga su _ . / ; \codec per didascalie e argomenti mobili)
\DeclareUrlCommand\code{\urlstyle{tt}}
\newcommand{\codec}[1]{\texttt{\detokenize{#1}}}
% Endpoint: monospazio, a capo sulle barre
\DeclareUrlCommand\ep{\urlstyle{tt}\def\UrlBreaks{\do\/}}


% ---------------------------------------------------------------------
%  Diagrammi di sequenza
% ---------------------------------------------------------------------
\tikzset{
  arr/.style={-{Stealth[length=2.2mm]}},
  comp/.style={draw, minimum width=2.2cm, minimum height=0.8cm,
                font=\sffamily\small, align=center, fill=white,
                text depth=0.3ex},
  db/.style={draw, cylinder, shape border rotate=90, aspect=0.22,
             minimum width=1.7cm, minimum height=0.95cm,
             font=\sffamily\scriptsize, align=center},
  state/.style={draw, rounded corners=6pt, minimum width=3cm,
                minimum height=0.9cm, font=\sffamily\small, align=center}
}
% \lifeline{x}{nome}{y finale}
\newcommand{\lifeline}[3]{%
  \node[draw, fill=white, minimum width=2.1cm, minimum height=0.65cm,
        font=\sffamily\footnotesize, text depth=0.25ex] at (#1,0) {#2};
  \draw[gray, densely dashed] (#1,-0.33) -- (#1,#3);}
% \msg{x da}{x a}{y}{etichetta}
\newcommand{\msg}[4]{%
  \draw[arr] (#1,#3) -- node[above, font=\ttfamily\scriptsize, align=center,
        inner sep=1.5pt, fill=white]{#4} (#2,#3);}
% \rmsg{x da}{x a}{y}{etichetta}: risposta
\newcommand{\rmsg}[4]{%
  \draw[arr, densely dashed] (#1,#3) -- node[above, font=\sffamily\scriptsize,
        align=center, inner sep=1.5pt, fill=white]{#4} (#2,#3);}
% \selfmsg{x}{y}{etichetta}
\newcommand{\selfmsg}[3]{%
  \draw[arr] (#1,#2) -- ++(0.3,0) -- ++(0,-0.3) -- ++(-0.3,0);
  \node[anchor=west, font=\sffamily\scriptsize, align=left]
       at ($(#1,#2)+(0.35,-0.15)$) {#3};}
% \frag{x1}{y1}{x2}{y2}{tipo}{condizione}
\newcommand{\frag}[6]{%
  \draw (#1,#2) rectangle (#3,#4);
  \node[draw, fill=white, anchor=north west, font=\sffamily\scriptsize\bfseries,
        inner sep=2pt] at (#1,#2) {#5};
  \node[anchor=north west, font=\sffamily\scriptsize] at ($(#1,#2)+(0.75,0)$) {[#6]};}
% \fragsep{x1}{x2}{y}{condizione}
\newcommand{\fragsep}[4]{%
  \draw[dashed] (#1,#3) -- (#2,#3);
  \node[anchor=north west, font=\sffamily\scriptsize] at (#1,#3) {[#4]};}

% ---------------------------------------------------------------------
%  Dati del documento
% ---------------------------------------------------------------------
\newcommand{\versione}{1.1}
\newcommand{\datadoc}{9 novembre 2026 (prima versione: 15 ottobre 2026)}
\newcommand{\gruppo}{Nome del gruppo e membri da inserire}

\hypersetup{
  pdftitle={Monster Tamers — Architettura del backend},
  pdfsubject={Advanced Software Engineering (Lab), Università di Pisa},
}

% =====================================================================
\begin{document}
% =====================================================================

% ---------------------------------------------------------------------
%  Frontespizio
% ---------------------------------------------------------------------
\hypersetup{pageanchor=false}
\begin{titlepage}
  \sffamily
  \vspace*{1cm}
  {\large Università di Pisa}\\[0.2em]
  {\normalsize Advanced Software Engineering — Laboratorio\quad·\quad A.A. 2026/27}

  \vspace{0.6cm}
  \rule{\linewidth}{0.8pt}

  \vspace{3.2cm}
  {\fontsize{34}{40}\selectfont\bfseries Monster Tamers}\\[0.8em]
  {\LARGE Architettura del backend\\[0.2em] e decisioni di progetto}

  \vspace{1.2cm}
  {\large Documento di progetto}

  \vfill
  \begin{tabular}{@{}ll@{}}
    \toprule
    Gruppo          & \gruppo \\
    Versione        & \versione \\
    Data            & \datadoc \\
    Stato           & Decisioni approvate dal gruppo; parametri di gioco da bilanciare \\
    Docente         & Alessandro Bocci \\
    \bottomrule
  \end{tabular}
  \vspace{1cm}
\end{titlepage}
\hypersetup{pageanchor=true}

% ---------------------------------------------------------------------
%  Sommario
% ---------------------------------------------------------------------
\section*{Sommario}

Questo documento descrive l'architettura del backend di \emph{Monster Tamers}, il gioco di collezione di mostri realizzato come progetto del laboratorio di Advanced Software Engineering, e motiva le decisioni prese dal gruppo.

Il backend è un'applicazione distribuita composta da \textbf{sei servizi Quarkus} (\code{auth}, \code{players}, \code{monsters}, \code{teams}, \code{encounters}, \code{history}), ciascuno con il proprio database PostgreSQL, esposti attraverso un unico \textbf{gateway Nginx}. I servizi comunicano solo tramite \textbf{REST sincrono}; l'autenticazione usa \textbf{JWT} firmati dal servizio \code{auth} e verificati localmente da tutti gli altri.

Il gioco è strutturato come un \textbf{roguelike a run}: il giocatore sceglie tre mostri iniziali, affronta encounter interattivi a turni in otto stage di difficoltà crescente e vince la run battendo il boss finale. Una sconfitta azzera la run; account, storico e run passate restano consultabili.

Il documento copre, nell'ordine: requisiti e vincoli (Sezione~\ref{sec:requisiti}), regole di gioco (Sezione~\ref{sec:gioco}), architettura generale (Sezione~\ref{sec:architettura}), dettaglio dei servizi con API e dati (Sezione~\ref{sec:servizi}), flussi principali (Sezione~\ref{sec:flussi}), consistenza dei dati (Sezione~\ref{sec:dati}), sicurezza (Sezione~\ref{sec:sicurezza}), deployment e test (Sezione~\ref{sec:deploy}), registro delle decisioni con le alternative scartate (Sezione~\ref{sec:decisioni}), modifiche introdotte dal piano di implementazione (Sezione~\ref{sec:aggiornamento}) e punti aperti (Sezione~\ref{sec:aperti}).

\tableofcontents
\clearpage

% =====================================================================
\section{Introduzione}
% =====================================================================

\subsection{Scopo del documento}

Il documento ha tre obiettivi:
\begin{enumerate}
  \item descrivere \textbf{cosa} fa il backend: componenti, responsabilità, API e dati;
  \item spiegare \textbf{perché} è fatto così, riportando per ogni decisione le alternative considerate e le conseguenze;
  \item servire da riferimento comune al gruppo durante lo sviluppo e da base per la consegna dell'architettura, per il report finale e per la discussione orale.
\end{enumerate}

\subsection{Contesto}

\emph{Monster Tamers} è un gioco di collezione di mostri. Il ciclo di gioco minimo richiesto dal docente è \emph{Prepara} $\rightarrow$ \emph{Affronta} $\rightarrow$ \emph{Esito} $\rightarrow$ \emph{Progresso}: il giocatore raccoglie mostri, forma un team, affronta un \emph{encounter} contro un mostro selvatico e, in base all'esito, fa progressi e ottiene nuovi mostri.

Il focus del progetto è il backend, che deve gestire collezione globale e collezioni personali, selezione del team, encounter e storico persistente, il tutto tramite un'interfaccia RESTful. Il frontend è minimale e non contiene logica di business.

\subsection{Documenti di riferimento}

\begin{tabularx}{\linewidth}{@{}L{6.4cm}Y@{}}
  \toprule
  Documento & Contenuto \\
  \midrule
  \codec{Project Kickoff.pdf} & Presentazione del progetto (07/10/2026) \\
  \code{Monster_Tamers_Backlog.xlsx} & 28 user stories con priorità \\
  \code{Monster_Tamers_Riepilogo.pdf} & Riepilogo in italiano di consegna e requisiti \\
  \code{area_giocatori_mostri.md} & Decisioni su giocatori e mostri (D1–D13, G1–G9) e acceptance criteria \\
  \code{architettura_backend.md} & Versione sintetica di questo documento \\
  \code{backend/monsters/seed/} & Catalogo dei mostri, abilità ed elementi (JSON) \\
  \bottomrule
\end{tabularx}

\subsection{Convenzioni}

\begin{itemize}
  \item Le decisioni sono identificate da un prefisso: \textbf{A} per l'architettura, \textbf{R} per le regole di gioco. Le decisioni \textbf{D} (dati) e \textbf{G} (catalogo) sono quelle del documento dell'area giocatori e mostri.
  \item I valori marcati come \emph{iniziali} (formule, soglie, livelli) sono parametri di bilanciamento: la struttura è decisa, i numeri possono cambiare dopo le simulazioni.
  \item Gli endpoint pubblici hanno il prefisso \ep{/api} e passano dal gateway; quelli con prefisso \ep{/internal} sono raggiungibili solo tra servizi.
  \item Nomi di servizi, tabelle e campi sono in inglese, come nel codice; il testo è in italiano.
\end{itemize}

\subsection{Glossario}

\begin{tabularx}{\linewidth}{@{}L{3.2cm}Y@{}}
  \toprule
  Termine & Significato \\
  \midrule
  Specie & Un mostro del catalogo globale (100 specie), identificato da \code{monster_id}. \\
  Istanza & Un mostro posseduto da un giocatore, identificato da \code{owned_id}, con livello e XP propri. Della stessa specie possono esistere più istanze. \\
  Run & Una partita completa: dalla scelta degli starter alla vittoria finale, alla sconfitta o alla resa. \\
  Stage & Livello di difficoltà della run, da 1 a 8. \\
  Template & Definizione statica di un encounter: stage, elemento, livelli e rarità dei selvatici possibili. \\
  Selvatico & Il mostro avversario di un encounter, estratto dal catalogo all'avvio. \\
  Boss & Encounter speciale di fine stage con un selvatico prefissato e più forte. \\
  Battaglia & Esecuzione concreta di un encounter: stato in corso, turni, esito. \\
  Staffetta & Formato in cui i mostri del team combattono uno alla volta, nell'ordine del team. \\
  \bottomrule
\end{tabularx}

% =====================================================================
\section{Requisiti e vincoli}\label{sec:requisiti}
% =====================================================================

\subsection{Requisiti tecnici}

\begin{tabularx}{\linewidth}{@{}L{3.6cm}YL{2.2cm}@{}}
  \toprule
  Requisito & Come viene soddisfatto & Sezione \\
  \midrule
  Applicazione distribuita & Sei servizi indipendenti in container separati, ognuno con il proprio database & \ref{sec:architettura} \\
  Interfaccia REST & Tutte le operazioni del giocatore sono endpoint REST documentati con OpenAPI & \ref{sec:servizi} \\
  Uno o più database & Un'istanza PostgreSQL per servizio; storico degli encounter persistente & \ref{sec:dati} \\
  Meccanismi di sicurezza & HTTPS, JWT, bcrypt, autorizzazione per risorsa, endpoint interni protetti & \ref{sec:sicurezza} \\
  Test unitari e di performance & JUnit, test di API con RestAssured, Locust sul gateway & \ref{sec:deploy} \\
  Analisi di sicurezza & Misure elencate per minaccia; analisi completa per la consegna finale & \ref{sec:sicurezza} \\
  Nessuna logica nel frontend & Regole, calcoli ed estrazioni casuali solo nel backend & \ref{sec:gioco} \\
  \bottomrule
\end{tabularx}

\subsection{User stories}

Il backlog contiene 28 user stories: 18 per il primo prototipo, 2 aggiuntive per la consegna finale (US-05 e US-25, entrambe sulla sicurezza) e 8 opzionali. Le scelte di questo documento coprono anche alcune opzionali senza costi aggiuntivi rilevanti: US-04 (nome e avatar), US-07 (scelta del primo mostro, tramite gli starter), US-13 (sviluppo dei mostri, tramite livelli e XP), US-21 (fuga) e US-22 (scelta delle azioni, tramite la battaglia interattiva). Non sono coperte US-08 (fazioni), US-26 e US-27 (amici e sfide). La corrispondenza completa tra user stories, servizi ed endpoint è in Appendice~\ref{app:tracciabilita}.

\subsection{Vincoli di progetto}

\begin{itemize}
  \item Gruppo di 3–4 persone; ogni membro deve saper discutere individualmente la propria parte all'orale.
  \item Consegna finale il 15 dicembre; prototipo opzionale il 10 novembre con le 18 user stories della prima consegna.
  \item Principio KISS richiesto dal docente: la complessità va aggiunta solo dove porta un beneficio misurabile.
  \item Framework scelto dal gruppo: Quarkus (Java).
\end{itemize}

% =====================================================================
\section{Il gioco}\label{sec:gioco}
% =====================================================================

\subsection{Visione d'insieme}

Il gioco è un \textbf{roguelike a run} (R1). Una run inizia con la scelta di tre mostri iniziali e si sviluppa in \textbf{otto stage} di difficoltà crescente (R2). In ogni stage il giocatore affronta encounter normali, che può ripetere, e un boss. Battere il boss dello stage 8 vince la run; perdere un encounter la fa terminare con un reset completo (R14).

Il ciclo di gioco richiesto dal docente si ritrova così:

\begin{tabularx}{\linewidth}{@{}L{2.4cm}Y@{}}
  \toprule
  Fase & In Monster Tamers \\
  \midrule
  Prepara & Scelta degli starter, consultazione della collezione, creazione e modifica dei team \\
  Affronta & Battaglia interattiva a turni contro un mostro selvatico \\
  Esito & Vittoria (il selvatico entra in collezione, XP al team), sconfitta (fine della run) o fuga \\
  Progresso & Storico degli encounter, livelli dei mostri, avanzamento di stage, run passate \\
  \bottomrule
\end{tabularx}

\begin{figure}[htbp]
  \centering
  \begin{adjustbox}{max width=\linewidth}
  \begin{tikzpicture}
    \node[circle, fill=black, inner sep=2.5pt] (start) at (-3.7,0) {};
    \node[state] (s1) at (0,0) {Starter\\da scegliere};
    \node[state] (s2) at (6.2,0) {Run attiva\\(stage 1–8)};
    \node[state] (s3) at (12.2,1.4) {Run vinta};
    \node[state] (s4) at (12.2,-1.4) {Run persa\\o abbandonata};

    \draw[arr] (start) -- node[above, font=\sffamily\scriptsize]{registrazione} (s1);
    \draw[arr] (s1) -- node[above, font=\sffamily\scriptsize]{scelta dei 3 starter} (s2);
    \draw[arr] ([xshift=-4mm]s2.north) arc[start angle=180, end angle=0, radius=0.4]
      node[pos=0.5, above, font=\sffamily\scriptsize, align=center]{encounter vinto,\\fuga, scadenza};
    \draw[arr] (s2.east) -- node[above, sloped, font=\sffamily\scriptsize]{boss stage 8 battuto} (s3.west);
    \draw[arr] (s2.east) -- node[below, sloped, font=\sffamily\scriptsize]{sconfitta o resa} (s4.west);
    \draw[arr] (s3.north) -- ++(0,0.6) -| node[pos=0.25, above, font=\sffamily\scriptsize]{nuova run richiesta dal giocatore, con reset} (s1.north);
    \draw[arr] (s4.south) -- ++(0,-0.6) -| node[pos=0.25, below, font=\sffamily\scriptsize]{reset automatico} (s1.south);
  \end{tikzpicture}
  \end{adjustbox}
  \caption{Ciclo di vita di una run dal punto di vista del giocatore.}
  \label{fig:run}
\end{figure}

\subsection{Giocatore}

Account e profilo sono separati (D1). L'\textbf{account} contiene solo le credenziali (username e hash della password) e vive nel servizio \code{auth}. Il \textbf{profilo} contiene i dati di gioco: nome visualizzato, avatar, run corrente e pool degli starter, e vive nel servizio \code{players}. I due condividono lo stesso identificativo \code{player_id}.

I progressi non sono duplicati in un unico campo: si ricavano dalla collezione (\code{monsters}), dalla run corrente (\code{encounters}) e dallo storico (\code{history}).

\subsection{Mostri}

Il catalogo contiene \textbf{100 specie} (D9), prodotte dalla pipeline in \code{tools/monster-gen/}. Ogni specie ha:
\begin{itemize}
  \item un'\textbf{immagine} pixel art 64×64 servita come file statico (D8);
  \item un \textbf{elemento} fra sette elementi concettuali;
  \item una \textbf{rarità}: \code{common}, \code{rare} o \code{legendary};
  \item quattro \textbf{statistiche} base intere 1–100 (\code{hp}, \code{attack}, \code{defense}, \code{speed}) con un tasso di \textbf{crescita} per livello (D5, D13);
  \item una \textbf{abilità} speciale scelta da un catalogo condiviso di 28 abilità (D12).
\end{itemize}

\begin{table}[htbp]
  \centering
  \small
  \begin{tabular}{@{}lrrrr@{}}
    \toprule
    Elemento & Common & Rare & Legendary & Totale \\
    \midrule
    Silence & 10 & 4 & 1 & 15 \\
    Echo    & 10 & 4 & 1 & 15 \\
    Memory  & 10 & 5 & 1 & 16 \\
    Time    &  9 & 4 & 2 & 15 \\
    Chaos   &  1 & 4 & 3 &  8 \\
    Gravity & 10 & 4 & 1 & 15 \\
    Dream   & 10 & 5 & 1 & 16 \\
    \midrule
    Totale  & 60 & 30 & 10 & 100 \\
    \bottomrule
  \end{tabular}
  \caption{Distribuzione del catalogo per elemento e rarità. La somma delle statistiche base cresce con la rarità: 148–174 per i common, 206–233 per i rare, 270–288 per i legendary.}
\end{table}

Gli elementi formano un \textbf{ciclo} (Figura~\ref{fig:elementi}): ogni elemento è forte contro il successivo, debole contro il precedente e neutro con gli altri quattro. Chaos è un elemento speciale con pochi mostri, quasi tutti rari o leggendari.

\begin{figure}[htbp]
  \centering
  \begin{tikzpicture}
    \foreach \n [count=\i from 0] in {Silence, Echo, Memory, Time, Chaos, Gravity, Dream} {
      \node[draw, rounded corners=3pt, minimum width=1.9cm, minimum height=0.65cm,
            font=\sffamily\small, fill=white] (e\i) at ({90-\i*360/7}:2.7cm) {\n};
    }
    \foreach \i in {0,...,6} {
      \pgfmathtruncatemacro\j{mod(\i+1,7)}
      \draw[arr] (e\i) -- (e\j);
    }
    \node[font=\sffamily\scriptsize, align=center] at (0,0) {A $\rightarrow$ B:\\A è forte contro B};
  \end{tikzpicture}
  \caption{Ciclo degli elementi (\codec{elements.json}). Un attacco forte ha moltiplicatore 1,5, uno debole 0,67.}
  \label{fig:elementi}
\end{figure}

Le \textbf{abilità} sono descritte da tre campi: tipo di effetto (\code{effect_type}), intensità (\code{effect_value}) e condizione di attivazione (\code{trigger}). La logica che le applica è solo nel backend. Il catalogo completo è in Appendice~\ref{app:abilita}.

\subsection{Collezione e istanze}

Ogni mostro posseduto è un'\textbf{istanza} con il proprio \code{owned_id}, livello ed esperienza; più istanze della stessa specie sono ammesse (D6). Ogni istanza appartiene a una run (\code{run_id}) e ha uno stato: \code{active} se appartiene alla run corrente, \code{archived} se appartiene a una run conclusa (R15). La collezione mostrata al giocatore è quella della run corrente; le collezioni delle run passate restano consultabili.

Le statistiche effettive dipendono dal livello e sono calcolate solo dal servizio \code{monsters}:
\[
  s_{\text{eff}} = s_{\text{base}} + g_s \cdot (L - 1)
\]
dove $s_{\text{base}}$ è la statistica base della specie, $g_s$ il suo tasso di crescita e $L$ il livello dell'istanza (massimo 50).

\subsection{Starter}

Alla creazione dell'account, e all'inizio di ogni nuova run, il backend estrae a caso un \textbf{pool di 10 specie common} e lo salva nel profilo, così ricaricare la pagina non cambia l'offerta. Il giocatore sceglie \textbf{3 specie distinte} dal pool; le tre istanze entrano in collezione a livello 1 e la run si apre allo stage 1. La scelta si fa una volta per run.

\subsection{Team}

Prima di un encounter il giocatore prepara un team con le istanze della propria collezione. Le regole sono volutamente poche (R8):
\begin{itemize}
  \item da \textbf{1 a 3} membri, numero che coincide con gli starter e permette di avere subito un team completo;
  \item ogni membro è un'istanza posseduta dal giocatore, attiva e appartenente alla run corrente;
  \item la stessa istanza non può comparire due volte, mentre più istanze della stessa specie sono ammesse;
  \item l'\textbf{ordine} dei membri conta: è l'ordine di entrata in battaglia.
\end{itemize}
Il team si può modificare liberamente tra un encounter e l'altro (R9). Il giocatore può salvare più team per run (valore iniziale: massimo 5).

\subsection{Encounter e stage}

Gli encounter disponibili sono definiti da \textbf{template statici} caricati come seed (R12). Ogni stage ha tre template normali e un boss. Un template descrive le caratteristiche dell'encounter (stage, elemento prevalente, intervallo di livello, pesi delle rarità) ma non il mostro esatto: il selvatico viene estratto dal catalogo all'avvio della battaglia, rispettando quelle caratteristiche. Il giocatore vede quindi cosa aspettarsi senza conoscere l'avversario preciso.

\begin{itemize}
  \item Sono disponibili i template degli stage già raggiunti, così il giocatore può tornare negli stage precedenti per fare esperienza.
  \item Il boss dello stage corrente si sblocca dopo \textbf{2 vittorie} nello stage (valore iniziale).
  \item Il boss ha una specie prefissata: \code{rare} negli stage 1–3, \code{legendary} di elementi diversi negli stage 4–7, \textbf{Omnivoid} (chaos, legendary) nello stage 8.
  \item Ogni boss si può battere \textbf{una volta per run} (R13): la vittoria dà il boss come ricompensa e fa salire lo stage.
\end{itemize}

\begin{table}[htbp]
  \centering
  \small
  \begin{tabular}{@{}lll@{}}
    \toprule
    Selvatico & Livello (stage $s$) & Esempio stage 8 \\
    \midrule
    Normale & casuale tra $3s-2$ e $3s$ & 22–24 \\
    Boss    & $3s+2$                    & 26 \\
    \bottomrule
  \end{tabular}
  \caption{Livelli dei selvatici (valori iniziali).}
\end{table}

\subsection{Battaglia}

La battaglia è \textbf{interattiva a turni} (R3) nel formato \textbf{a staffetta} (R4): combatte un solo mostro del team alla volta, nell'ordine del team, contro il selvatico, che conserva i danni ricevuti per tutta la battaglia. Tutti i mostri iniziano con gli HP pieni (R10).

\subsubsection*{Azioni}

A ogni turno il giocatore sceglie una delle tre azioni (R5):

\begin{tabularx}{\linewidth}{@{}L{2.6cm}Y@{}}
  \toprule
  Azione & Effetto \\
  \midrule
  Attacco & Il mostro attivo attacca il selvatico. \\
  Cambio  & Un altro membro del team, non KO, prende il posto del mostro attivo. Il cambio consuma il turno. \\
  Fuga    & Tenta di chiudere l'encounter. Riesce con una probabilità che dipende dalla velocità (R7); se fallisce, il turno è perso. \\
  \bottomrule
\end{tabularx}

Non esiste un'azione di cattura: la ricompensa è legata alla vittoria (R6).

\subsubsection*{Svolgimento del turno}

Ogni richiesta del giocatore risolve un \textbf{round completo}: l'azione del giocatore e quella del selvatico, che attacca sempre.
\begin{enumerate}
  \item Cambio e fuga si risolvono per primi.
  \item Gli attacchi si risolvono in ordine di velocità effettiva; a parità decide il caso.
  \item Se il mostro attivo va KO entra automaticamente il successivo ancora in piedi.
  \item La battaglia termina con la vittoria (selvatico KO), con la sconfitta (tutto il team KO) o con una fuga riuscita.
\end{enumerate}

\subsubsection*{Formule}

Valori iniziali, da bilanciare con uno script di simulazione:
\begin{align*}
  HP_{\max} &= 2\,hp_{\text{eff}} + 10
    && \times 1{,}5 \text{ selvatico},\ \times 2{,}5 \text{ boss} \\
  D &= \max\!\left(1,\ \operatorname{round}\!\left(\frac{atk^2}{atk + def}\cdot 0{,}6 \cdot m_{\text{el}} \cdot v\right)\right)
    && v \sim U[0{,}9;\ 1{,}1] \\
  m_{\text{el}} &\in \{1{,}5;\ 1;\ 0{,}67\}
    && \\
  p_{\text{fuga}} &= \min\!\left(0{,}95,\ \frac{spd_{a}}{spd_{a} + spd_{w}} + 0{,}25\right)
    && \\
  XP &= 10 \cdot L_{w} \cdot f_r
    && \\
  XP_{L \to L+1} &= 15\,L
    && L \le 50
\end{align*}

Legenda: $atk$ e $def$ sono attacco dell'attaccante e difesa del difensore; $m_{\text{el}}$ vale 1,5 se l'attacco è forte, 1 se neutro, 0,67 se debole; $spd_a$ e $spd_w$ sono le velocità del mostro attivo e del selvatico; $L_w$ è il livello del selvatico e $f_r$ vale 1, 1,5 o 2 per common, rare e legendary.

Gli HP del selvatico sono moltiplicati perché combatte da solo contro fino a tre avversari.

\subsubsection*{Abilità in battaglia}

Le abilità di entrambi i contendenti si applicano secondo il loro \code{trigger}. Gli effetti \emph{continui} (riduzione del danno, aumento del danno, schivata) valgono finché la condizione è vera; gli effetti \emph{puntuali} (\code{heal}, \code{revive}) scattano una sola volta per battaglia.

\begin{tabularx}{\linewidth}{@{}L{3cm}Y@{}}
  \toprule
  Trigger & Quando si applica \\
  \midrule
  \code{always}      & Per tutta la battaglia \\
  \code{first_turn}  & Alla prima azione del mostro \\
  \code{hp_below_50} & Con HP sotto il 50\% (effetti puntuali: la prima volta che scendono sotto) \\
  \code{hp_above_50} & Con HP sopra il 50\% \\
  \code{on_hit}      & Quando il mostro subisce un attacco \\
  \code{on_attack}   & Quando il mostro attacca \\
  \code{vs_faster}   & Contro un avversario più veloce \\
  \code{once}        & Una volta per battaglia, all'evento previsto dall'effetto \\
  \bottomrule
\end{tabularx}

\subsubsection*{Esiti}

\begin{tabularx}{\linewidth}{@{}L{3cm}YL{4.4cm}@{}}
  \toprule
  Esito & Ricompensa & Effetto sulla run \\
  \midrule
  \code{victory}     & Il selvatico entra in collezione al proprio livello; XP a tutti i membri del team (R11) & Continua; se boss, stage $+1$ o run vinta \\
  \code{defeat}      & Nessuna & Run persa: reset completo \\
  \code{fled}        & Nessuna & Continua \\
  \code{expired}     & Nessuna (battaglia inattiva da 30 minuti) & Continua \\
  \code{interrupted} & Nessuna (riavvio del servizio) & Continua \\
  \bottomrule
\end{tabularx}

Il mostro ottenuto entra in collezione al \textbf{livello del selvatico}, non a livello 1: in una struttura a run un mostro di livello 1 trovato allo stage 6 sarebbe inutilizzabile e non premierebbe la vittoria.

\subsection{Run: fine, reset e resa}

\begin{itemize}
  \item \textbf{Sconfitta} (R14): la run diventa \code{lost} e parte il reset. Le istanze della run passano ad \code{archived}, i team della run vengono eliminati, il profilo riceve una nuova run con un nuovo pool di starter. Account, storico e run passate non vengono toccati.
  \item \textbf{Resa} (R17): il giocatore può abbandonare la run in qualsiasi momento fuori da una battaglia; la run diventa \code{abandoned} con lo stesso reset.
  \item \textbf{Vittoria finale} (R16): battuto il boss dello stage 8, la run diventa \code{won}. Il giocatore inizia una nuova run quando vuole; l'avvio esegue lo stesso reset.
  \item \textbf{Run passate} (R15): restano consultabili run, collezione archiviata e storico, così il giocatore può capire dove ha sbagliato.
\end{itemize}

% =====================================================================
\section{Architettura}\label{sec:architettura}
% =====================================================================

\subsection{Stile architetturale}

Il backend segue uno stile a \textbf{microservizi}: ogni componente con uno scopo chiaro diventa un servizio Quarkus indipendente, con il proprio database e il proprio ciclo di build e deploy. La scelta risponde direttamente al requisito di applicazione distribuita e permette a ogni membro del gruppo di essere responsabile di uno o due servizi, cosa utile anche per la valutazione individuale.

La consegna dell'architettura chiede uno schizzo dei componenti; i componenti individuati qui coincidono con i servizi, quindi lo schizzo e l'implementazione restano allineati.

\subsection{Vista dei componenti}

\begin{figure}[htbp]
  \centering
  \begin{adjustbox}{max width=\linewidth}
  \begin{tikzpicture}
    \node[comp, minimum width=3.4cm] (client) at (6.5,4.4) {Client\\browser e frontend};
    \node[comp, minimum width=4.6cm] (gw) at (6.5,2.6) {Nginx — API gateway\\HTTPS, routing, rate limiting};
    \node[comp, dashed] (fe) at (11.9,2.6) {Frontend\\statico};
    \foreach \s [count=\i from 0] in {auth, players, monsters, teams, encounters, history} {
      \node[comp] (s\i) at (\i*2.6,0.6) {\texttt{\s}};
      \node[db] (d\i) at (\i*2.6,-1.3) {\s-db};
      \draw (s\i) -- (d\i);
      \draw[arr] (gw.south) -- (s\i.north);
    }
    \draw[arr] (client) -- node[right, font=\sffamily\scriptsize]{HTTPS} (gw);
    \draw[dashed] (gw) -- (fe);
    \node[draw, densely dotted, inner sep=0.35cm, fit=(s0)(s5)(d0)(d5)] (net) {};
    \node[anchor=north east, font=\sffamily\scriptsize] at (net.south east) {rete Docker privata};
  \end{tikzpicture}
  \end{adjustbox}
  \caption{Vista dei componenti. Solo il gateway espone una porta; servizi e database sono raggiungibili solo sulla rete interna.}
  \label{fig:componenti}
\end{figure}

\subsection{Responsabilità dei servizi}

{\small
\begin{tabularx}{\linewidth}{@{}L{2.1cm}YL{4.1cm}L{2.2cm}@{}}
  \toprule
  Servizio & Responsabilità & Dati & User stories \\
  \midrule
  \code{auth} & Registrazione, login, token dei giocatori e dei servizi & \code{accounts}, \code{refresh_tokens} & US-01, 02, 05 \\
  \code{players} & Profilo, pool degli starter, run corrente & \code{players} & US-03, 04, 06, 07 \\
  \code{monsters} & Catalogo, abilità, elementi, immagini; collezione per run, livelli ed XP & \code{monsters}, \code{abilities}, \code{elements}, \code{owned_monsters} & US-09, 10, 11, 12, 13 \\
  \code{teams} & Creazione, modifica, validazione dei team & \code{teams}, \code{team_slots} & US-14, 15, 16 \\
  \code{encounters} & Catalogo degli encounter, run, battaglie, orchestrazione di fine run & \code{encounter_templates}, \code{runs}; battaglie in memoria & US-17–22 \\
  \code{history} & Storico persistente e integrità dei salvataggi & \code{encounter_records} & US-23, 24, 25 \\
  \bottomrule
\end{tabularx}}

\subsection{Gateway}

Nginx è l'unico punto di ingresso (A3). Svolge quattro compiti: instradare le richieste verso il servizio competente, terminare HTTPS, limitare la frequenza delle richieste sugli endpoint di autenticazione e servire il frontend statico. Non contiene logica applicativa né verifica i token: questo compito resta ai servizi.

\begin{tabularx}{\linewidth}{@{}YL{2.6cm}@{}}
  \toprule
  Percorso pubblico & Servizio \\
  \midrule
  \ep{/api/auth/*} & \code{auth} \\
  \ep{/api/players/*} & \code{players} \\
  \ep{/api/monsters/*}, \ep{/api/abilities/*}, \ep{/api/elements}, \ep{/api/collection/*}, \ep{/images/*} & \code{monsters} \\
  \ep{/api/teams/*} & \code{teams} \\
  \ep{/api/encounters/*}, \ep{/api/runs/*}, \ep{/api/battles/*} & \code{encounters} \\
  \ep{/api/history/*} & \code{history} \\
  \ep{/} & frontend \\
  \bottomrule
\end{tabularx}

I percorsi \ep{/internal/*} non sono mai instradati.

\subsection{Comunicazione tra servizi}

I servizi comunicano solo tramite \textbf{REST sincrono} con i REST client tipizzati di Quarkus (A5). Le chiamate interne sono protette dalle annotazioni di SmallRye Fault Tolerance (\code{@Timeout}, \code{@Retry}) e sono \textbf{idempotenti}: ogni chiamata che modifica dati porta una chiave (\code{run_id} o \code{encounter_id}) che permette di ripeterla senza effetti doppi.

\begin{figure}[htbp]
  \centering
  \begin{tikzpicture}[lbl/.style={font=\sffamily\scriptsize, fill=white, inner sep=1.5pt}]
    \node[comp] (auth) at (0,2.8) {\code{auth}};
    \node[comp] (pl) at (5,2.8) {\code{players}};
    \node[comp] (te) at (0,0) {\code{teams}};
    \node[comp] (en) at (5,0) {\code{encounters}};
    \node[comp] (mo) at (10.8,0) {\code{monsters}};
    \node[comp] (hi) at (5,-2.6) {\code{history}};

    \draw[arr] (auth) -- node[lbl]{crea profilo} (pl);
    \draw[arr] (pl) -- node[lbl, sloped]{starter} (mo);
    \draw[arr] ([xshift=-4mm]pl.south) -- node[lbl, left=2pt]{apre run} ([xshift=-4mm]en.north);
    \draw[arr] ([xshift=4mm]en.north) -- node[lbl, right=2pt]{nuova run} ([xshift=4mm]pl.south);
    \draw[arr] ([yshift=1.5mm]en.west) -- node[lbl, above=1pt]{team} ([yshift=1.5mm]te.east);
    \draw[arr] ([yshift=-1.5mm]en.west) -- node[lbl, below=1pt]{reset} ([yshift=-1.5mm]te.east);
    \draw[arr] (en) -- node[lbl, above=2pt, align=center, fill=none]{statistiche, ricompense,\\XP, archivio} (mo);
    \draw[arr] (en) -- node[lbl]{salva encounter} (hi);
    \draw[arr] (te.south) to[bend right=26] node[lbl, pos=0.75]{verifica possesso} (mo.south);
  \end{tikzpicture}
  \caption{Dipendenze tra servizi. Ogni freccia è una chiamata a un endpoint \codec{/internal/*} con token di servizio. \codec{history} non chiama nessuno; \codec{encounters} orchestra la fine delle battaglie.}
  \label{fig:dipendenze}
\end{figure}

Il grafo delle dipendenze non ha cicli tra chiamate sincrone nello stesso flusso: \code{players} chiama \code{encounters} solo all'apertura della run, \code{encounters} chiama \code{players} solo alla chiusura.

\subsection{Autenticazione e autorizzazione}

Il servizio \code{auth} firma i token con una chiave privata RSA (A6); gli altri servizi montano solo la chiave pubblica e verificano i token localmente, senza chiamare \code{auth} a ogni richiesta.

\begin{itemize}
  \item \textbf{Access token}: JWT RS256 con \code{sub} = \code{player_id} e ruolo \code{player}, durata 15 minuti.
  \item \textbf{Refresh token}: valore casuale, salvato in \code{auth} solo come hash, durata 7 giorni, revocabile con il logout.
  \item \textbf{Token di servizio}: JWT con ruolo \code{service}, ottenuto da ciascun servizio presso \code{auth} con credenziali proprie (\emph{client credentials}) e riusato fino alla scadenza.
  \item \textbf{Autorizzazione per risorsa}: il \code{player_id} si legge sempre dal token, mai da parametri del client. Un giocatore non può quindi indicare le risorse di un altro.
\end{itemize}

Gli endpoint \ep{/internal/*} hanno due livelli di protezione (A7): non sono raggiungibili dall'esterno perché il gateway non li instrada, e richiedono comunque un token con ruolo \code{service}.

\subsection{Stack tecnologico}

\begin{tabularx}{\linewidth}{@{}L{3.2cm}L{4.4cm}Y@{}}
  \toprule
  Ambito & Tecnologia & Motivo \\
  \midrule
  Linguaggio e framework & Java 21, Quarkus (LTS 3.x) & Scelta del gruppo; avvio rapido, Dev Services, estensioni pronte \\
  API REST & Quarkus REST + Jackson, SmallRye OpenAPI & Specifica OpenAPI generata dal codice \\
  Persistenza & Hibernate ORM con Panache, Flyway & Modello più documentato; migrazioni versionate \\
  Database & PostgreSQL & Relazionale, supporto JSONB per log e snapshot \\
  Sicurezza & SmallRye JWT, bcrypt, Hibernate Validator & Verifica locale dei token; hash robusto; validazione dichiarativa \\
  Chiamate interne & REST Client, SmallRye Fault Tolerance & Client tipizzati; timeout e retry dichiarativi \\
  Gateway & Nginx & Configurazione dichiarativa, nessun codice \\
  Build e deploy & Maven, Docker, docker compose & Un comando per buildare, uno per avviare \\
  Test & JUnit 5, RestAssured, WireMock, Locust & Unit, API, integrazione, performance \\
  \bottomrule
\end{tabularx}

% =====================================================================
\section{Servizi}\label{sec:servizi}
% =====================================================================

Per ogni servizio sono riportati endpoint e dati. Tutti gli endpoint pubblici, tranne registrazione e login, richiedono l'intestazione \codec{Authorization: Bearer} con un access token valido. Le chiavi di idempotenza degli endpoint interni sono indicate nella descrizione.

% ---------------------------------------------------------------------
\subsection{\texttt{auth}}

{\small\begin{tabularx}{\linewidth}{@{}L{1.3cm}L{6.1cm}Y@{}}
  \toprule
  Metodo & Endpoint & Descrizione \\
  \midrule
  POST & \ep{/api/auth/register} & Crea l'account da \code{username} e \code{password} e restituisce il \code{player_id}. Crea il profilo su \code{players}; se la chiamata fallisce, l'account viene annullato. \\
  POST & \ep{/api/auth/login} & Restituisce access e refresh token. Con credenziali errate restituisce un errore generico, senza indicare quale campo è sbagliato. \\
  POST & \ep{/api/auth/refresh} & Emette un nuovo access token da un refresh token valido. \\
  POST & \ep{/api/auth/logout} & Revoca il refresh token. \\
  POST & \ep{/internal/auth/service-token} & Emette un token di servizio da \code{client_id} e \code{client_secret}. \\
  \bottomrule
\end{tabularx}}

{\small\begin{tabularx}{\linewidth}{@{}L{4.0cm}L{2.9cm}Y@{}}
  \toprule
  Campo & Tipo & Note \\
  \midrule
  \multicolumn{3}{@{}l}{\textit{Tabella} \code{accounts}} \\
  \code{player_id} & UUID, PK & Stesso id del profilo \\
  \code{username} & varchar(20), unique & 3–20 caratteri \\
  \code{password_hash} & varchar & bcrypt con salt; la password in chiaro non viene mai salvata \\
  \code{created_at}, \code{last_login} & timestamp & \\
  \midrule
  \multicolumn{3}{@{}l}{\textit{Tabella} \code{refresh_tokens}} \\
  \code{token_hash} & char(64), PK & SHA-256 del token \\
  \code{player_id} & UUID & \\
  \code{expires_at} & timestamp & \\
  \code{revoked} & boolean & \\
  \bottomrule
\end{tabularx}}

La chiave privata di firma esiste solo in \code{auth}. Le credenziali dei servizi sono passate come variabili d'ambiente.

% ---------------------------------------------------------------------
\subsection{\texttt{players}}

{\small\begin{tabularx}{\linewidth}{@{}L{1.3cm}L{6.1cm}Y@{}}
  \toprule
  Metodo & Endpoint & Descrizione \\
  \midrule
  GET & \ep{/api/players/me} & Profilo, run corrente e stato degli starter. \\
  PATCH & \ep{/api/players/me} & Modifica nome visualizzato e avatar (US-04). \\
  GET & \ep{/api/players/me/starters} & Pool di 10 common della run corrente. \\
  POST & \ep{/api/players/me/starters} & Sceglie 3 starter distinti dal pool; li aggiunge in \code{monsters} e apre la run in \code{encounters}. \\
  POST & \ep{/internal/players} & Crea il profilo; idempotente su \code{player_id}. \\
  POST & \ep{/internal/players/{pid}/new-run} & Assegna una nuova run e un nuovo pool. Riceve la run conclusa e non fa nulla se non è quella corrente: per questo è idempotente. \\
  \bottomrule
\end{tabularx}}

{\small\begin{tabularx}{\linewidth}{@{}L{4.0cm}L{2.9cm}Y@{}}
  \toprule
  Campo & Tipo & Note \\
  \midrule
  \multicolumn{3}{@{}l}{\textit{Tabella} \code{players}} \\
  \code{player_id} & UUID, PK & \\
  \code{display_name} & varchar(30) & Inizialmente uguale allo username \\
  \code{avatar} & varchar, nullable & \\
  \code{current_run_id} & UUID & Generato dal servizio a ogni nuova run \\
  \code{starter_pool} & int[10] & Specie common estratte per la run corrente \\
  \code{starters_chosen} & boolean & \code{true} dopo la scelta; torna \code{false} al reset \\
  \code{created_at} & timestamp & \\
  \bottomrule
\end{tabularx}}

Lo stage non è nel profilo ma nella run (\code{encounters}), perché si azzera a ogni run.

% ---------------------------------------------------------------------
\subsection{\texttt{monsters}}

{\small\begin{tabularx}{\linewidth}{@{}L{1.3cm}L{6.1cm}Y@{}}
  \toprule
  Metodo & Endpoint & Descrizione \\
  \midrule
  GET & \ep{/api/monsters} & Catalogo paginato, filtrabile per \code{element} e \code{rarity} (US-10). \\
  GET & \ep{/api/monsters/{id}} & Dettaglio della specie con abilità (US-11). \\
  GET & \ep{/api/abilities}, \ep{/api/abilities/{id}} & Catalogo delle abilità. \\
  GET & \ep{/api/elements} & Ciclo degli elementi. \\
  GET & \ep{/api/collection} & Collezione della run corrente; con \code{run_id} quella archiviata di una run passata, solo del giocatore stesso (US-09). \\
  GET & \ep{/api/collection/{owned_id}} & Istanza con livello, XP, statistiche effettive, provenienza (US-11). \\
  GET & \ep{/images/monsters/{NNN}.png} & Sprite statico. \\
  POST & \ep{/internal/collections/{pid}/starters} & Aggiunge i 3 starter; idempotente su \code{run_id}. \\
  POST & \ep{/internal/collections/{pid}/monsters} & Aggiunge il mostro vinto al livello indicato; idempotente su \code{encounter_id} (US-12). \\
  POST & \ep{/internal/collections/{pid}/xp} & Assegna XP alle istanze e calcola i level up; idempotente su \code{encounter_id}. \\
  GET & \ep{/internal/collections/{pid}/instances} & Istanze richieste con statistiche effettive e abilità. \\
  GET & \ep{/internal/monsters/{id}/stats} & Statistiche di una specie a un dato livello, per il selvatico. \\
  POST & \ep{/internal/collections/{pid}/archive} & Archivia le istanze di una run; idempotente su \code{run_id}. \\
  \bottomrule
\end{tabularx}}

{\small\begin{tabularx}{\linewidth}{@{}L{4.0cm}L{2.9cm}Y@{}}
  \toprule
  Campo & Tipo & Note \\
  \midrule
  \multicolumn{3}{@{}l}{\textit{Tabelle di catalogo} \code{monsters}, \code{abilities}, \code{elements}} \\
  \multicolumn{3}{@{}p{\linewidth}@{}}{Struttura identica ai file di seed; caricate da una migrazione Flyway e mai modificate a runtime. \code{monsters.ability_id} è una chiave esterna reale, perché le tabelle sono nello stesso database.} \\
  \midrule
  \multicolumn{3}{@{}l}{\textit{Tabella} \code{owned_monsters}} \\
  \code{owned_id} & UUID, PK & Referenziato da team e storico \\
  \code{player_id} & UUID & Riferimento logico a \code{players} \\
  \code{run_id} & UUID & Run di appartenenza \\
  \code{monster_id} & int, FK & Specie \\
  \code{level} & smallint & 1–50 \\
  \code{xp} & int & XP accumulata nel livello corrente \\
  \code{source} & enum & \code{starter}, \code{encounter}, \code{boss} \\
  \code{encounter_id} & UUID, unique, nullable & Encounter di provenienza; garantisce l'idempotenza \\
  \code{status} & enum & \code{active}, \code{archived} \\
  \code{obtained_at} & timestamp & \\
  \midrule
  \multicolumn{3}{@{}l}{\textit{Tabella} \code{xp_grants}} \\
  \code{encounter_id} & UUID, PK & XP già assegnata per quell'encounter \\
  \code{granted_at} & timestamp & \\
  \bottomrule
\end{tabularx}}

La formula delle statistiche effettive esiste solo in questo servizio; \code{teams} ed \code{encounters} la usano tramite gli endpoint interni.

% ---------------------------------------------------------------------
\subsection{\texttt{teams}}

{\small\begin{tabularx}{\linewidth}{@{}L{1.3cm}L{6.1cm}Y@{}}
  \toprule
  Metodo & Endpoint & Descrizione \\
  \midrule
  GET & \ep{/api/teams} & Team della run corrente. \\
  POST & \ep{/api/teams} & Crea un team con nome e membri ordinati, dopo averlo validato (US-14). \\
  GET & \ep{/api/teams/{id}} & Dettaglio del team. \\
  PUT & \ep{/api/teams/{id}} & Modifica membri e ordine (US-15). \\
  DELETE & \ep{/api/teams/{id}} & Elimina il team. \\
  POST & \ep{/api/teams/validate} & Valida una composizione senza salvarla; restituisce l'elenco degli errori (US-16). \\
  GET & \ep{/internal/teams/{id}} & Team rivalidato, con le istanze, per l'avvio di una battaglia. \\
  POST & \ep{/internal/teams/{pid}/reset} & Elimina i team di una run; idempotente su \code{run_id}. \\
  \bottomrule
\end{tabularx}}

{\small\begin{tabularx}{\linewidth}{@{}L{4.0cm}L{2.9cm}Y@{}}
  \toprule
  Campo & Tipo & Note \\
  \midrule
  \multicolumn{3}{@{}l}{\textit{Tabella} \code{teams}} \\
  \code{team_id} & UUID, PK & \\
  \code{player_id}, \code{run_id} & UUID & \\
  \code{name} & varchar(30) & \\
  \code{created_at}, \code{updated_at} & timestamp & \\
  \midrule
  \multicolumn{3}{@{}l}{\textit{Tabella} \code{team_slots}} \\
  \code{team_id} & UUID, FK & PK insieme a \code{position} \\
  \code{position} & smallint & 1–3, ordine di entrata in battaglia \\
  \code{owned_id} & UUID & Unico all'interno del team \\
  \bottomrule
\end{tabularx}}

La validazione chiede a \code{monsters} le istanze indicate e controlla: numero di membri tra 1 e 3, possesso, stato \code{active}, appartenenza alla run corrente, assenza di istanze ripetute. Un team valido al salvataggio viene rivalidato all'avvio di ogni battaglia.

% ---------------------------------------------------------------------
\subsection{\texttt{encounters}}

{\small\begin{tabularx}{\linewidth}{@{}L{1.3cm}L{6.1cm}Y@{}}
  \toprule
  Metodo & Endpoint & Descrizione \\
  \midrule
  GET & \ep{/api/encounters} & Encounter disponibili: template degli stage raggiunti e boss dello stage corrente, se sbloccato (US-17). \\
  GET & \ep{/api/encounters/{template_id}} & Caratteristiche del template. \\
  GET & \ep{/api/runs} & Run del giocatore, correnti e passate, con esito e stage raggiunto. \\
  GET & \ep{/api/runs/current} & Stage, vittorie nello stage, boss battuti. \\
  GET & \ep{/api/runs/{run_id}} & Dettaglio di una run. \\
  POST & \ep{/api/runs/current/abandon} & Resa: chiude la run ed esegue il reset. \\
  POST & \ep{/api/runs} & Nuova run dopo una run vinta; esegue il reset. \\
  POST & \ep{/api/battles} & Avvia una battaglia da template e team; 409 se ne esiste già una (US-18). \\
  GET & \ep{/api/battles/current} & Stato della battaglia: mostro attivo, HP, turno, selvatico. \\
  POST & \ep{/api/battles/current/actions} & Esegue un'azione (\code{attack}, \code{switch}, \code{flee}) e restituisce l'esito del round (US-19, US-21, US-22). Alla fine della battaglia restituisce esito e ricompense (US-20). \\
  POST & \ep{/internal/runs} & Apre una run allo stage 1; idempotente su \code{run_id}. \\
  \bottomrule
\end{tabularx}}

{\small\begin{tabularx}{\linewidth}{@{}L{4.0cm}L{2.9cm}Y@{}}
  \toprule
  Campo & Tipo & Note \\
  \midrule
  \multicolumn{3}{@{}l}{\textit{Tabella} \code{encounter_templates} (seed)} \\
  \code{template_id} & int, PK & \\
  \code{stage} & smallint & 1–8 \\
  \code{name}, \code{description} & varchar, text & \\
  \code{element} & varchar, nullable & Elemento prevalente; \code{null} = qualsiasi \\
  \code{rarity_weights} & JSONB & Probabilità di estrazione per rarità \\
  \code{level_min}, \code{level_max} & smallint & \\
  \code{is_boss} & boolean & \\
  \code{boss_monster_id} & int, nullable & Specie del boss \\
  \midrule
  \multicolumn{3}{@{}l}{\textit{Tabella} \code{runs}} \\
  \code{run_id} & UUID, PK & Generato da \code{players} \\
  \code{player_id} & UUID & Al massimo una run \code{active} per giocatore (indice parziale) \\
  \code{status} & enum & \code{active}, \code{won}, \code{lost}, \code{abandoned} \\
  \code{stage} & smallint & 1–8 \\
  \code{wins_in_stage} & smallint & Per lo sblocco del boss \\
  \code{bosses_defeated} & int[] & Stage il cui boss è stato battuto \\
  \code{active_battle_id} & UUID, nullable & Battaglia in corso \\
  \code{reset_pending} & boolean & Reset da completare (Sezione~\ref{sec:dati}) \\
  \code{version} & int & Lock ottimistico (\code{@Version}) \\
  \code{started_at}, \code{ended_at} & timestamp & \\
  \bottomrule
\end{tabularx}}

\subsubsection*{Battaglie in memoria}

Lo stato delle battaglie in corso è tenuto in memoria (R19), in una mappa indicizzata per giocatore: snapshot del team, HP correnti, mostro attivo, effetti attivi, turni giocati, seed del generatore casuale e istante dell'ultima azione. Le azioni di una stessa battaglia sono elaborate in modo sequenziale.

La run, invece, è persistita: all'avvio di una battaglia il suo id viene salvato in \code{runs.active_battle_id}. Questo permette di gestire due casi:
\begin{itemize}
  \item \textbf{Scadenza}: un job eseguito ogni minuto chiude come \code{expired} le battaglie inattive da più di 30 minuti (R18).
  \item \textbf{Riavvio}: all'avvio il servizio cerca le run con \code{active_battle_id} valorizzato ma senza battaglia in memoria e le chiude come \code{interrupted}.
\end{itemize}
In entrambi i casi l'encounter viene registrato nello storico e la run continua.

% ---------------------------------------------------------------------
\subsection{\texttt{history}}

{\small\begin{tabularx}{\linewidth}{@{}L{1.3cm}L{6.1cm}Y@{}}
  \toprule
  Metodo & Endpoint & Descrizione \\
  \midrule
  GET & \ep{/api/history} & Encounter passati del giocatore, paginati e filtrabili per \code{run_id} ed esito (US-23). \\
  GET & \ep{/api/history/{encounter_id}} & Dettaglio completo con team, selvatico e log dei turni (US-24). \\
  POST & \ep{/internal/history} & Registra un encounter concluso; idempotente su \code{encounter_id}. \\
  \bottomrule
\end{tabularx}}

{\small\begin{tabularx}{\linewidth}{@{}L{4.0cm}L{2.9cm}Y@{}}
  \toprule
  Campo & Tipo & Note \\
  \midrule
  \multicolumn{3}{@{}l}{\textit{Tabella} \code{encounter_records}} \\
  \code{encounter_id} & UUID, PK & \\
  \code{player_id}, \code{run_id} & UUID & \\
  \code{template_id}, \code{stage} & int, smallint & \\
  \code{wild} & JSONB & Specie, livello e statistiche del selvatico \\
  \code{team_snapshot} & JSONB & Istanze con livello e statistiche al momento della battaglia \\
  \code{outcome} & enum & \code{victory}, \code{defeat}, \code{fled}, \code{expired}, \code{interrupted} \\
  \code{reward_owned_id} & UUID, nullable & Istanza ottenuta \\
  \code{xp_gained} & int & \\
  \code{turns} & JSONB & Log dei turni: azione, danni, abilità attivate \\
  \code{rng_seed} & bigint & Permette di rigiocare la battaglia \\
  \code{started_at}, \code{ended_at} & timestamp & \\
  \code{prev_hash}, \code{hash} & char(64) & Catena di integrità (Sezione~\ref{sec:sicurezza}) \\
  \bottomrule
\end{tabularx}}

Lo storico salva uno snapshot completo (R20): il dettaglio di un encounter passato resta leggibile anche dopo che le istanze sono state archiviate o hanno cambiato livello.

% =====================================================================
\section{Flussi principali}\label{sec:flussi}
% =====================================================================

Nei diagrammi il gateway è omesso: tutte le richieste del client passano da Nginx. Le frecce continue sono richieste, quelle tratteggiate risposte.

\subsection{Registrazione e scelta degli starter}

\begin{figure}[H]
  \centering
  \begin{adjustbox}{max width=\linewidth}
  \begin{tikzpicture}
    \lifeline{0}{Client}{-11.6}
    \lifeline{4.0}{auth}{-11.6}
    \lifeline{8.0}{players}{-11.6}
    \lifeline{12.0}{monsters}{-11.6}
    \lifeline{16.0}{encounters}{-11.6}

    \msg{0}{4.0}{-1.0}{POST /api/auth/register}
    \selfmsg{4.0}{-1.4}{salva account\\(hash bcrypt)}
    \msg{4.0}{8.0}{-2.5}{POST /internal/players}
    \selfmsg{8.0}{-2.9}{genera run\_id e\\pool di 10 common}
    \rmsg{8.0}{4.0}{-4.0}{201}
    \rmsg{4.0}{0}{-4.6}{201, player\_id}
    \msg{0}{4.0}{-5.5}{POST /api/auth/login}
    \rmsg{4.0}{0}{-6.1}{access e refresh token}
    \msg{0}{8.0}{-7.0}{POST /api/players/me/starters}
    \selfmsg{8.0}{-7.4}{valida la scelta}
    \msg{8.0}{12.0}{-8.6}{POST /internal/collections/\\\{pid\}/starters}
    \rmsg{12.0}{8.0}{-9.2}{201}
    \msg{8.0}{16.0}{-10.0}{POST /internal/runs}
    \rmsg{16.0}{8.0}{-10.6}{201, stage 1}
    \rmsg{8.0}{0}{-11.2}{200, starter scelti}
  \end{tikzpicture}
  \end{adjustbox}
  \caption{Registrazione, login e scelta degli starter.}
\end{figure}

Se la creazione del profilo fallisce, \code{auth} annulla l'account appena creato e restituisce un errore: non restano account senza profilo. La scelta degli starter segna \code{starters_chosen} solo dopo che entrambe le chiamate interne sono riuscite; essendo idempotenti, una richiesta ripetuta dal client completa il lavoro senza duplicati.

\subsection{Avvio e svolgimento di una battaglia}

\begin{figure}[H]
  \centering
  \begin{adjustbox}{max width=\linewidth}
  \begin{tikzpicture}
    \lifeline{0}{Client}{-10.4}
    \lifeline{4.6}{encounters}{-10.4}
    \lifeline{9.2}{teams}{-10.4}
    \lifeline{13.8}{monsters}{-10.4}

    \msg{0}{4.6}{-1.1}{POST /api/battles\\template\_id, team\_id}
    \msg{4.6}{9.2}{-1.8}{GET /internal/teams/\{id\}}
    \msg{9.2}{13.8}{-2.6}{GET /internal/collections/\\\{pid\}/instances}
    \rmsg{13.8}{9.2}{-3.2}{istanze e statistiche}
    \rmsg{9.2}{4.6}{-3.8}{team valido}
    \msg{4.6}{13.8}{-4.6}{GET /internal/monsters/\{id\}/stats}
    \rmsg{13.8}{4.6}{-5.2}{statistiche del selvatico}
    \selfmsg{4.6}{-5.6}{crea la battaglia in memoria,\\salva active\_battle\_id}
    \rmsg{4.6}{0}{-6.6}{stato iniziale}

    \frag{-0.7}{-7.0}{7.6}{-10.0}{loop}{ogni turno}
    \msg{0}{4.6}{-8.3}{POST /api/battles/\\current/actions}
    \selfmsg{4.6}{-8.7}{risolve il round}
    \rmsg{4.6}{0}{-9.6}{esito del round e stato}
  \end{tikzpicture}
  \end{adjustbox}
  \caption{Avvio di una battaglia e ciclo dei turni.}
\end{figure}

Durante i turni \code{encounters} non chiama altri servizi: tutto quello che serve è nello snapshot preso all'avvio. Il costo di una mossa è quindi una sola richiesta HTTP, che non dipende dalla disponibilità degli altri servizi.

\subsection{Fine della battaglia e reset della run}

\begin{figure}[H]
  \centering
  \begin{adjustbox}{max width=\linewidth}
  \begin{tikzpicture}
    \lifeline{0}{encounters}{-8.8}
    \lifeline{3.8}{history}{-8.8}
    \lifeline{7.6}{monsters}{-8.8}
    \lifeline{11.4}{teams}{-8.8}
    \lifeline{15.2}{players}{-8.8}

    \frag{-0.7}{-0.7}{16.1}{-8.4}{alt}{vittoria}
    \msg{0}{7.6}{-1.7}{POST /internal/collections/\{pid\}/monsters}
    \msg{0}{7.6}{-2.4}{POST /internal/collections/\{pid\}/xp}
    \selfmsg{0}{-2.8}{aggiorna la run:\\vittorie, stage}
    \msg{0}{3.8}{-4.0}{POST /internal/history}

    \fragsep{-0.7}{16.1}{-4.5}{sconfitta o resa}
    \selfmsg{0}{-5.2}{run: lost o abandoned}
    \msg{0}{3.8}{-6.1}{POST /internal/history}
    \msg{0}{7.6}{-6.7}{POST /internal/collections/\{pid\}/archive}
    \msg{0}{11.4}{-7.3}{POST /internal/teams/\{pid\}/reset}
    \msg{0}{15.2}{-7.9}{POST /internal/players/\{pid\}/new-run}
  \end{tikzpicture}
  \end{adjustbox}
  \caption{Fine di una battaglia: ricompense in caso di vittoria, reset della run in caso di sconfitta o resa. Le risposte sono omesse.}
  \label{fig:fine}
\end{figure}

Il reset è una \textbf{saga orchestrata} da \code{encounters} (Sezione~\ref{sec:dati}). La resa segue lo stesso ramo, senza il passo sullo storico se non c'era una battaglia in corso.

% =====================================================================
\section{Dati e consistenza}\label{sec:dati}
% =====================================================================

\subsection{Un database per servizio}

Ogni servizio ha il proprio container PostgreSQL e le proprie credenziali (A4); nessun servizio può leggere i dati di un altro. Ne derivano due conseguenze:
\begin{itemize}
  \item i riferimenti tra dati di servizi diversi (per esempio \code{owned_monsters.player_id} o \code{team_slots.owned_id}) sono \textbf{logici}: si salva l'identificativo, ma il database non ne garantisce l'esistenza. Il controllo è fatto dal servizio, tramite le chiamate interne;
  \item non esistono transazioni che coinvolgono più servizi: la consistenza tra servizi va gestita esplicitamente.
\end{itemize}

\subsection{Consistenza tra servizi}

Le operazioni che toccano più servizi sono tre: registrazione, scelta degli starter e fine della battaglia. Tutte seguono le stesse regole:
\begin{enumerate}
  \item ogni chiamata interna che modifica dati è \textbf{idempotente} grazie a una chiave naturale dell'operazione (\code{player_id}, \code{run_id}, \code{encounter_id}), protetta da un vincolo \code{unique} nel database del servizio chiamato;
  \item in caso di errore temporaneo le chiamate sono ripetute automaticamente, con \code{@Retry} e \code{@Timeout};
  \item lo stato locale del chiamante viene aggiornato prima delle chiamate, così da sapere sempre cosa resta da fare.
\end{enumerate}

Per il reset della run, \code{encounters} imposta \code{status} a \code{lost} (o \code{abandoned}) e \code{reset_pending} a \code{true}, poi esegue in ordine archiviazione della collezione, reset dei team e nuova run sul profilo. Quando tutti i passi riescono, \code{reset_pending} torna \code{false}. Se un passo fallisce in modo definitivo, il flag resta attivo e il reset riprende alla richiesta successiva del giocatore: grazie all'idempotenza, i passi già eseguiti non hanno effetti doppi.

\subsection{Migrazioni e seed}

Lo schema di ogni database è definito da migrazioni Flyway versionate nel repository del servizio. I dati statici (catalogo dei mostri, abilità, elementi, template degli encounter) sono caricati da migrazioni dedicate e non sono modificabili tramite API. I file di seed dei mostri sono generati dalla pipeline \code{tools/monster-gen/} e non vanno modificati a mano.

\subsection{Immagini}

Gli sprite dei mostri sono file statici serviti da \code{monsters}; nel database c'è solo il percorso (\code{image_url}). Il frontend li mostra ingranditi senza sfocatura, con la proprietà CSS \codec{image-rendering: pixelated}.

% =====================================================================
\section{Sicurezza}\label{sec:sicurezza}
% =====================================================================

Le misure sono progettate fin da ora, anche se US-05 e US-25 appartengono alla consegna finale: introdurle più tardi richiederebbe di modificare tutti i servizi.

\begin{xltabular}{\linewidth}{@{}L{4.3cm}YL{2.6cm}@{}}
  \toprule
  Minaccia & Contromisura & Dove \\
  \midrule
  \endhead
  Intercettazione del traffico & HTTPS terminato sul gateway (certificato self-signed in sviluppo) & Nginx \\
  Furto del database delle credenziali & Password salvate solo come hash bcrypt con salt; refresh token salvati come hash & \code{auth} \\
  Furto di un token & Access token di 15 minuti; refresh token revocabile & \code{auth} \\
  Accesso ai dati di un altro giocatore & \code{player_id} letto solo dal token; ogni risorsa verificata contro il proprietario & tutti \\
  Chiamata diretta a endpoint interni & \ep{/internal/*} non instradati; ruolo \code{service} obbligatorio & Nginx, tutti \\
  Brute force su login e registrazione & Rate limiting per indirizzo IP & Nginx \\
  Injection e input malformati & Validazione dichiarativa (Hibernate Validator); query parametrizzate tramite ORM & tutti \\
  Imbrogli sul gioco & Regole, estrazioni casuali e ricompense solo nel backend; nessun endpoint per aggiungere mostri o XP & \code{encounters}, \code{monsters} \\
  Manomissione dello storico & Tabella append-only (l'utente applicativo non ha \code{UPDATE} né \code{DELETE}); catena di hash SHA-256 per giocatore; seed salvato per rigiocare e verificare una battaglia & \code{history} \\
  Perdita dei salvataggi & Dati persistenti su volumi Docker; le run concluse vengono archiviate, non cancellate & tutti \\
  Segreti nel repository & Credenziali e chiavi solo da variabili d'ambiente & docker compose \\
  \bottomrule
\end{xltabular}

La catena di hash funziona così: ogni record contiene l'hash del record precedente dello stesso giocatore e il proprio hash, calcolato su \code{prev_hash} e sul contenuto del record. Modificare o cancellare un record rompe la catena da quel punto in poi, e la manomissione diventa rilevabile.

L'analisi di sicurezza completa, con gli strumenti che verranno presentati nel corso, è prevista per la consegna finale.

% =====================================================================
\section{Deployment e test}\label{sec:deploy}
% =====================================================================

\subsection{Deployment}

L'intero sistema si avvia con docker compose: Nginx, sei servizi Quarkus e sei container PostgreSQL, collegati da una rete Docker privata. Solo Nginx espone una porta all'esterno.

\begin{table}[htbp]
  \centering
  \small
  \begin{tabularx}{\linewidth}{@{}L{7.2cm}Y@{}}
    \toprule
    Percorso & Contenuto \\
    \midrule
    \code{backend/pom.xml} & pom aggregatore: versione di Quarkus, Java 21, moduli \\
    \code{backend/<servizio>/} & progetto Maven del servizio, \code{Dockerfile} \\
    \code{backend/<servizio>/src/main/resources/db/migration/} & migrazioni Flyway \\
    \code{gateway/nginx.conf} & configurazione del gateway \\
    \code{frontend/} & frontend statico \\
    \code{docker-compose.yml} & topologia completa \\
    \bottomrule
  \end{tabularx}
  \caption{Organizzazione del repository.}
\end{table}

I servizi non condividono codice (A9): ognuno definisce i propri DTO. La duplicazione è limitata e viene accettata in cambio dell'indipendenza: una modifica a un servizio non obbliga a ricompilare gli altri.

\subsection{Strategia di test}

\begin{tabularx}{\linewidth}{@{}L{3.4cm}YL{3.4cm}@{}}
  \toprule
  Livello & Oggetto & Strumenti \\
  \midrule
  Unitario & Formule, regole di validazione del team, motore di battaglia con seed fisso (risultati deterministici) & JUnit 5 \\
  API del servizio & Endpoint, codici di risposta, autorizzazione & \code{@QuarkusTest}, RestAssured, Dev Services \\
  Integrazione & Chiamate verso altri servizi, retry, idempotenza & WireMock, mock dei REST client \\
  End-to-end & Flusso completo tramite gateway & docker compose, RestAssured \\
  Performance & Login, catalogo, battaglie concorrenti & Locust \\
  \bottomrule
\end{tabularx}

Il motore di battaglia usa un generatore casuale con seed: con lo stesso seed e le stesse azioni produce sempre lo stesso risultato. Questo rende i test deterministici e permette di verificare a posteriori una battaglia registrata nello storico.

% =====================================================================
\section{Registro delle decisioni}\label{sec:decisioni}
% =====================================================================

Per ogni decisione sono riportate la scelta, le alternative considerate e la motivazione.

\subsection{Architettura}

\begin{xltabular}{\linewidth}{@{}L{0.8cm}L{3.3cm}L{4.4cm}Y@{}}
  \toprule
  ID & Decisione & Alternative scartate & Motivazione e conseguenze \\
  \midrule
  \endhead
  A1 & Quarkus, Java 21 & — & Scelta del gruppo. Dev Services, estensioni pronte per REST, JWT, ORM, fault tolerance. \\
  A2 & 6 servizi & 7 (catalogo separato dalla collezione); 5 (storico dentro \code{encounters}); 4 accorpati & Confini chiari e un servizio per responsabilità; lavoro divisibile tra i membri. Costo: più container e chiamate interne. \\
  A3 & Gateway Nginx & Gateway Quarkus custom; Traefik; nessun gateway & Solo configurazione: routing, HTTPS, rate limiting. Senza gateway servirebbero CORS e HTTPS su ogni servizio. \\
  A4 & PostgreSQL, un container per servizio & Un PostgreSQL con database separati; PostgreSQL + MongoDB per lo storico & Isolamento reale dei dati e dei guasti; una sola tecnologia di database. Costo: più container, comunque leggeri. \\
  A5 & Solo REST sincrono & REST + broker di eventi (RabbitMQ, Kafka) & KISS: un solo paradigma, test semplici. La resilienza si ottiene con retry e idempotenza invece che con un broker. \\
  A6 & JWT RS256 firmati da \code{auth} & Keycloak con OIDC; token opachi verificati da \code{auth} & Verifica locale senza chiamate ad \code{auth}; nessun collo di bottiglia. La revoca è gestita con token brevi e refresh token revocabili. \\
  A7 & \ep{/internal/*} non instradati + token di servizio & Solo isolamento di rete; mTLS & Doppia difesa a costo moderato. mTLS resta un'estensione possibile. \\
  A8 & Hibernate Panache + Flyway & Hibernate Reactive; JDBC o jOOQ & Modello più semplice e documentato; il vantaggio del reattivo non è visibile a questa scala. \\
  A9 & Maven, pom aggregatore, nessun codice condiviso & Modulo \code{common} condiviso; progetti indipendenti; Gradle & Un comando di build senza accoppiare i servizi tramite una libreria comune. \\
  \bottomrule
\end{xltabular}

\subsection{Regole di gioco}

\begin{xltabular}{\linewidth}{@{}L{0.8cm}L{3.3cm}L{4.4cm}Y@{}}
  \toprule
  ID & Decisione & Alternative scartate & Motivazione e conseguenze \\
  \midrule
  \endhead
  R1 & Roguelike a run & Progressione continua senza reset & Ogni run ha inizio e fine; la sconfitta ha un peso reale. Introduce l'entità run e il reset tra servizi. \\
  R2 & 8 stage con boss finale & Zone per elemento aperte; obiettivo di completare il catalogo & Obiettivo chiaro e difficoltà crescente, come nell'esempio del docente. \\
  R3 & Battaglia interattiva a turni & Simulata in una chiamata; simulata con tattica iniziale; formula a punteggio & Il giocatore prende decisioni a ogni turno; copre anche US-21 e US-22. Richiede stato tra richieste. \\
  R4 & 1v1 a staffetta & Tutto il team contro il selvatico & Regole di danno semplici; l'ordine del team diventa una scelta strategica. \\
  R5 & Attacco, cambio, fuga & Anche un'azione di cattura & La cattura è stata esclusa per coerenza con R6. \\
  R6 & Vittoria = selvatico in collezione & Ottenimento con probabilità per rarità; cattura come azione & Regola semplice; ogni vittoria fa crescere la collezione. \\
  R7 & Fuga con probabilità basata sulla velocità & Fuga sempre riuscita & La velocità conta anche fuori dall'attacco; la fuga ha un rischio. \\
  R8 & Team da 1 a 3, copie della stessa specie ammesse, nessun altro vincolo & Massimo 4 o 6; specie distinte e massimo un legendary; budget di rarità & Coincide con i 3 starter; regole minime e facili da validare. \\
  R9 & Team modificabile tra un encounter e l'altro & Team bloccato per tutta la run & Libertà strategica; i mostri appena ottenuti si usano subito. \\
  R10 & HP pieni a ogni encounter & HP persistenti nella run, con o senza cura & Ogni encounter è indipendente; stato della run più semplice. \\
  R11 & XP a tutto il team & Solo a chi è sceso in campo; nessun livello & Semplice e generoso; permette di far crescere i mostri in coda al team. \\
  R12 & Template statici, selvatico estratto & Encounter generati per giocatore a rotazione; selvatico fisso & Dati semplici, varietà nelle battaglie, nessuno stato per giocatore. \\
  R13 & Boss battibile una volta per run, dà il boss & Boss ripetibile; boss non ottenibile & Rende i legendary trofei della run e impedisce di accumularli. \\
  R14 & Sconfitta: reset completo & Nessuna penalità; perdita di XP; mostri a riposo; reset del solo stage; reset dello stage e dei mostri della run & Scelta del gruppo per un roguelike puro. Richiede il reset coordinato di quattro servizi. \\
  R15 & Istanze archiviate, non cancellate & Cancellazione & Le run passate restano consultabili; nessuna perdita di dati irreversibile. \\
  R16 & Run vinta, poi nuova run & Modalità libera dopo la vittoria & Ciclo coerente: ogni run ha un inizio e una fine. \\
  R17 & Resa ammessa & Nessuna resa & Utile con un pool di starter sfortunato; riusa la logica di reset. \\
  R18 & Una battaglia attiva, scadenza dopo 30 minuti & Nessuna scadenza; battaglie in parallelo & Evita battaglie abbandonate per sempre e conflitti sullo stesso mostro. \\
  R19 & Battaglia in memoria, run in PostgreSQL & Tutto in PostgreSQL; Redis; tutto in memoria & Battaglia veloce e semplice da gestire; la run sopravvive ai riavvii. Limite: battaglie perse al riavvio e una sola istanza di \code{encounters}. \\
  R20 & Storico con riepilogo e log dei turni & Solo riepilogo & Dettaglio completo per US-24 e base per la verifica di integrità. \\
  \bottomrule
\end{xltabular}

\subsection{Modifiche alle decisioni precedenti}

L'introduzione delle run modifica alcune decisioni del documento dell'area giocatori e mostri:

\begin{tabularx}{\linewidth}{@{}L{5.6cm}Y@{}}
  \toprule
  Prima & Ora \\
  \midrule
  \code{Players.stage} & Spostato in \code{runs.stage}, nel servizio \code{encounters} \\
  Starter scelti una sola volta (D7) & Una volta per run \\
  Mostro ottenuto a livello 1 (US-12) & Al livello del selvatico \\
  \code{OwnedMonsters} senza run & Aggiunti \code{run_id} e \code{status} \\
  Progressi sempre conservati (US-03) & Conservati finché la run non termina; le run passate restano consultabili \\
  \bottomrule
\end{tabularx}

% =====================================================================
\section{Piano di implementazione (aggiornamento del 9/11)}\label{sec:aggiornamento}
% =====================================================================

Il piano di implementazione (\code{docs/plan/}, 6 novembre 2026) ha modificato alcune scelte di questo documento. Dove le sezioni precedenti e questa sezione non coincidono, vale questa sezione.

\subsection{Modifiche rispetto a questo documento}

\begin{xltabular}{\linewidth}{@{}L{6.2cm}Y@{}}
  \toprule
  Questo documento (15/10) & Piano di implementazione (06/11) \\
  \midrule
  \endhead
  POST \ep{/api/auth/login} (JSON) & POST \ep{/api/auth/token}, endpoint OAuth2 (form-urlencoded, \code{grant_type=password}; \code{refresh_token} in P2) \\
  POST \ep{/internal/auth/service-token} & POST \ep{/internal/auth/token} con \code{grant_type=client_credentials} \\
  Nginx serve anche il frontend statico su \ep{/} & Il frontend gira fuori da compose (Vite); Nginx fa solo da proxy per l'API, con CORS \\
  HTTPS terminato su Nginx & HTTPS su ogni collegamento: client$\to$gateway, gateway$\to$servizio, servizio$\to$servizio, servizio$\to$PostgreSQL \\
  Violazione di ownership $\to$ 403 (AC di US-05) & $\to$ \textbf{404}, senza rivelare se la risorsa esiste \\
  Body di POST \ep{/internal/players} non specificato & \texttt{\{player\_id\}}; il nome visualizzato di default lo genera \code{players} \\
  Nessun endpoint interno per filtrare il catalogo & GET \texttt{/internal/monsters?rarity=\&element=} \\
  US-20 senza endpoint & Risposta finale dell'azione di battaglia + GET \ep{/api/history/{id}} \\
  Esiti: \code{victory}, \code{defeat}, \code{fled}, \code{expired}, \code{interrupted} & + \code{abandoned} (battaglia chiusa dall'abbandono della run) \\
  Catalogo solo con autenticazione (D11) & Vale per gli endpoint JSON; le immagini dei mostri sono pubbliche (P-16) \\
  \bottomrule
\end{xltabular}

\needspace{8\baselineskip}
\subsection{Registro delle decisioni del piano (P-01…P-21)}

\begin{xltabular}{\linewidth}{@{}L{0.9cm}L{3.6cm}Y@{}}
  \toprule
  ID & Tema & Decisione \\
  \midrule
  \endhead
  P-01 & Ambito & Prototipo P1 (10/11) e consegna finale P2 (15/12) in un unico piano \\
  P-02 & Lingua & Codice, API, OpenAPI e piano in inglese; i documenti italiani precedenti restano \\
  P-03 & Frontend & Vite + React (TypeScript), fuori da Docker Compose \\
  P-04 & Segreti & Generati al primo avvio dal container \code{init-secrets} (password dei database, certificati, chiavi JWT, segreti dei client); nulla di segreto nel repository \\
  P-05 & Login & Endpoint OAuth2 POST \ep{/api/auth/token}; i servizi usano \code{client_credentials} su \ep{/internal/auth/token} \\
  P-06 & TLS & HTTPS ovunque già in P1: una CA locale, un certificato per container \\
  P-07 & Porta del gateway & \code{https://localhost:8443} \\
  P-08 & Stage & Tutti gli 8 stage già in P1 \\
  P-09 & Abilità & Effetti semplici in P1, complessi in P2 \\
  P-10 & Fuga e abbandono & Entrambi in P1: la fuga è un'azione di battaglia, l'abbandono riusa il reset della sconfitta \\
  P-11 & Risorse altrui & Sempre 404, mai 403 \\
  P-12 & Bruno & Collection di component test per tutti e 6 i servizi, più integration ed e2e \\
  P-13 & Accesso per i test & \code{docker-compose.test.yml} espone i servizi su localhost solo per i test \\
  P-14 & Cifratura a riposo (P2) & Username cifrato (AES-GCM) con indice HMAC; refresh token salvati solo come hash SHA-256 \\
  P-15 & Strumenti di sicurezza (P2) & Semgrep (SAST) e Trivy (dipendenze e immagini); in CI dal 09/11 \\
  P-16 & Immagini & Gli sprite dei mostri sono pubblici: \texttt{<img src>} non può inviare un bearer token \\
  P-17 & Token nel frontend & \code{localStorage}, con CSP e nessun HTML grezzo renderizzato \\
  P-18 & CI in P1 & GitHub Actions con build e test da P1; scansioni di sicurezza in P2 \\
  P-19 & Nome visualizzato di default & \texttt{Tamer-<prime 6 cifre esadecimali del player\_id>}; lo username non viene copiato in \code{players} \\
  P-20 & Robustezza della battaglia in P1 & Recupero pigro: una battaglia sparita dalla memoria viene chiusa come \code{interrupted}, una inattiva come \code{expired}, al successivo accesso del giocatore; sweep periodico in P2 \\
  P-21 & Boss & I boss di stage non compaiono negli incontri selvatici: si ottengono solo battendoli \\
  \bottomrule
\end{xltabular}

\subsection{Altre scelte fatte durante l'implementazione}

\begin{itemize}
  \item Moltiplicatore degli HP dei boss $\times 1{,}5$ (vedi \code{docs/plan/07-encounters.md}, \emph{Tuned constants}).
  \item Outbox nel servizio \code{encounters}: se un servizio non risponde, ricompensa, XP e record dello storico restano in sospeso e vengono ritentati al successivo accesso del giocatore alla run (idempotenti per \code{encounter_id}).
  \item Gateway: invia ai servizi \texttt{Host: <nome del servizio>}, non l'header \code{Host} del client (CI-2).
  \item Immagini Docker: PostgreSQL derivato senza \code{gosu} (\code{infra/postgres}), Jackson 2.22.4 sopra il BOM di Quarkus (CI-2, vedi \code{docs/security/}).
\end{itemize}

% =====================================================================
\section{Limiti noti e punti aperti}\label{sec:aperti}
% =====================================================================

\subsection{Limiti noti}

\begin{itemize}
  \item \textbf{Battaglie in memoria} (R19): un riavvio di \code{encounters} interrompe le battaglie in corso e il servizio non può essere replicato su più istanze. Run, stage e collezione non si perdono. Il passaggio a uno stato persistente non cambierebbe le API.
  \item \textbf{Consistenza eventuale durante il reset}: tra la fine di una run e il completamento del reset, alcuni servizi possono mostrare ancora i dati della run conclusa per un breve intervallo.
  \item \textbf{Comunicazione sincrona}: se un servizio chiamato è irraggiungibile, l'operazione che lo coinvolge fallisce dopo i retry; le battaglie in corso non ne sono toccate.
\end{itemize}

\subsection{Punti aperti}

\begin{itemize}
  \item Bilanciamento di formule, livelli e soglie con uno script di simulazione.
  \item Definizione dei 32 template degli encounter e scelta dei boss degli stage 1–7.
  \item Dettaglio degli effetti \code{copy_ability}, \code{invert_advantage} e \code{random_boost}.
  \item Analisi di sicurezza completa per la consegna finale.
\end{itemize}

% =====================================================================
\appendix
\section{Tracciabilità delle user stories}\label{app:tracciabilita}
% =====================================================================

Priorità: \textbf{P} primo prototipo, \textbf{F} consegna finale, \textbf{O} opzionale.

\begin{xltabular}{\linewidth}{@{}L{1.3cm}L{4.2cm}L{0.6cm}L{2.1cm}Y@{}}
  \toprule
  ID & Titolo & Pr. & Servizio & Endpoint principali \\
  \midrule
  \endhead
  US-01 & Creare un account & P & \code{auth} & POST \ep{/api/auth/register} \\
  US-02 & Accedere (login) & P & \code{auth} & POST \ep{/api/auth/login} \\
  US-03 & Ritrovare i progressi salvati & P & \code{players} & GET \ep{/api/players/me}, GET \ep{/api/runs/current} \\
  US-04 & Personalizzare il profilo & O & \code{players} & PATCH \ep{/api/players/me} \\
  US-05 & Proteggere i dati dell'account & F & tutti & Sezione~\ref{sec:sicurezza} \\
  US-06 & Ricevere una collezione iniziale & P & \code{players} & GET, POST \ep{/api/players/me/starters} \\
  US-07 & Scegliere il primo mostro & O & \code{players} & come US-06 \\
  US-08 & Scegliere una fazione & O & — & non coperta \\
  US-09 & Vedere la collezione personale & P & \code{monsters} & GET \ep{/api/collection} \\
  US-10 & Sfogliare la collezione globale & P & \code{monsters} & GET \ep{/api/monsters} \\
  US-11 & Vedere i dettagli di un mostro & P & \code{monsters} & GET \ep{/api/monsters/{id}}, GET \ep{/api/collection/{owned_id}} \\
  US-12 & Ottenere un mostro guadagnato & P & \code{monsters} & interno, a fine battaglia \\
  US-13 & Sviluppare i mostri & O & \code{monsters} & livelli e XP, interno \\
  US-14 & Creare un team & P & \code{teams} & POST \ep{/api/teams} \\
  US-15 & Modificare un team & P & \code{teams} & PUT \ep{/api/teams/{id}} \\
  US-16 & Validare un team & P & \code{teams} & POST \ep{/api/teams/validate} \\
  US-17 & Sfogliare gli encounter & P & \code{encounters} & GET \ep{/api/encounters} \\
  US-18 & Avviare un encounter & P & \code{encounters} & POST \ep{/api/battles} \\
  US-19 & Risolvere un encounter & P & \code{encounters} & POST \ep{/api/battles/current/actions} \\
  US-20 & Dettagli del risultato & P & \code{encounters} & risposta dell'ultima azione \\
  US-21 & Abbandonare un encounter & O & \code{encounters} & azione \code{flee}; resa della run \\
  US-22 & Scegliere le azioni & O & \code{encounters} & azioni \code{attack}, \code{switch}, \code{flee} \\
  US-23 & Vedere lo storico & P & \code{history} & GET \ep{/api/history} \\
  US-24 & Dettagli di un encounter passato & P & \code{history} & GET \ep{/api/history/{encounter_id}} \\
  US-25 & Proteggere i salvataggi & F & \code{history}, tutti & Sezione~\ref{sec:sicurezza} \\
  US-26 & Gestire gli amici & O & — & non coperta \\
  US-27 & Sfidare un amico & O & — & non coperta \\
  US-28 & Usare un'interfaccia semplice & P & frontend & tutte le API pubbliche \\
  \bottomrule
\end{xltabular}

% =====================================================================
\section{Catalogo delle abilità}\label{app:abilita}
% =====================================================================

Valori da \code{backend/monsters/seed/abilities.json}. L'ambito indica se l'abilità è generica, legata a un elemento o esclusiva dei legendary.

\begin{xltabular}{\linewidth}{@{}rL{3.1cm}L{1.9cm}L{3.6cm}rY@{}}
  \toprule
  ID & Nome & Ambito & Effetto & Valore & Trigger \\
  \midrule
  \endhead
  1  & Thick Hide          & generica & \code{damage_reduction}    & 10  & \code{always} \\
  2  & Quick Reflexes      & generica & \code{dodge}               & 10  & \code{always} \\
  3  & Second Wind         & generica & \code{heal}                & 15  & \code{hp_below_50} \\
  4  & Opening Strike      & generica & \code{damage_boost}        & 20  & \code{first_turn} \\
  5  & Stubborn            & generica & \code{damage_reduction}    & 20  & \code{hp_below_50} \\
  6  & Momentum            & generica & \code{speed_boost}         & 15  & \code{always} \\
  7  & Spiteful Thorns     & generica & \code{reflect}             & 10  & \code{on_hit} \\
  8  & Lullaby             & dream    & \code{stun}                & 15  & \code{on_attack} \\
  9  & Lucid Shift         & dream    & \code{dodge}               & 20  & \code{hp_above_50} \\
  10 & Resonance           & echo     & \code{damage_boost}        & 15  & \code{always} \\
  11 & Feedback            & echo     & \code{reflect}             & 20  & \code{on_hit} \\
  12 & Rewind              & time     & \code{heal}                & 25  & \code{hp_below_50} \\
  13 & Haste               & time     & \code{speed_boost}         & 25  & \code{first_turn} \\
  14 & Heavy Body          & gravity  & \code{damage_reduction}    & 20  & \code{always} \\
  15 & Crushing Pull       & gravity  & \code{damage_boost}        & 25  & \code{vs_faster} \\
  16 & Muffle              & silence  & \code{weaken}              & 15  & \code{first_turn} \\
  17 & Unseen              & silence  & \code{dodge}               & 25  & \code{first_turn} \\
  18 & Recall              & memory   & \code{damage_boost}        & 20  & \code{hp_below_50} \\
  19 & Imprint             & memory   & \code{ignore_disadvantage} & 100 & \code{always} \\
  20 & Entropy             & chaos    & \code{random_boost}        & 30  & \code{always} \\
  21 & Wild Surge          & chaos    & \code{damage_boost}        & 40  & \code{hp_below_50} \\
  22 & Dreamweaver         & legendary & \code{stun}               & 35  & \code{on_attack} \\
  23 & Endless Chorus      & legendary & \code{reflect}            & 40  & \code{on_hit} \\
  24 & Hourglass Reversal  & legendary & \code{revive}             & 50  & \code{once} \\
  25 & Event Horizon       & legendary & \code{damage_reduction}   & 40  & \code{always} \\
  26 & Absolute Quiet      & legendary & \code{weaken}             & 35  & \code{always} \\
  27 & Eternal Archive     & legendary & \code{copy_ability}       & 100 & \code{first_turn} \\
  28 & Paradox Engine      & legendary & \code{invert_advantage}   & 100 & \code{always} \\
  \bottomrule
\end{xltabular}

\begin{tabularx}{\linewidth}{@{}L{3.6cm}Y@{}}
  \toprule
  Effetto & Significato in battaglia \\
  \midrule
  \code{damage_boost}        & Danno inflitto aumentato del valore (\%) \\
  \code{damage_reduction}    & Danno subito ridotto del valore (\%) \\
  \code{dodge}               & Probabilità (\%) di evitare un attacco \\
  \code{heal}                & Recupera il valore (\%) degli HP massimi, una volta \\
  \code{speed_boost}         & Velocità aumentata del valore (\%) \\
  \code{reflect}             & Restituisce all'attaccante il valore (\%) del danno subito \\
  \code{stun}                & Probabilità (\%) che l'avversario salti la prossima azione \\
  \code{weaken}              & Attacco dell'avversario ridotto del valore (\%) \\
  \code{revive}              & Una volta per battaglia, al KO torna con il valore (\%) degli HP \\
  \code{random_boost}        & A ogni azione una statistica casuale aumenta del valore (\%) per quell'azione \\
  \code{ignore_disadvantage} & Lo svantaggio di elemento non si applica \\
  \code{invert_advantage}    & Vantaggi e svantaggi di elemento si invertono per entrambi \\
  \code{copy_ability}        & Copia l'abilità dell'avversario per il resto della battaglia \\
  \bottomrule
\end{tabularx}

\end{document}
